% Target 3 replacement proof, 21 November 2026.
% This is a section fragment: the parent supplies its document preamble.
\section{An independently reconstructed G34 source}
\label{sec:HD-G}

\paragraph{Provenance and scope.}
The ten vectors and seventeen intervals below are candidate data from
Severn, \emph{Weil classes on abelian sixfolds and eightfolds of Weil type
for $\mathbb Q(\sqrt{-3})$}, version 1.5, 21 October 2026,
Hexagon 2610.00387v2, Appendix A. That manuscript is unrefereed.
No theorem of that manuscript is a premise of the arguments below.
The group and all parabolic monodromy operators are reconstructed from
these vectors and intervals by \texttt{g34\_certificate.py}; its quotient
basis is chosen by elimination, without using the basis or matrices printed
in the manuscript. The proof of differential rank uses the resulting
$K$-linear span, the several-variable nilpotent orbit theorem, and the
one-variable limiting mixed Hodge structure theorem. The cycle argument
below supplies its own proof of the local sheaf assertion needed in
Schoen's ramified construction and of the transfer from Brauer induction.
The result is an independent proof verification of the requested case;
no priority claim is made for conclusions already stated by Severn.

\paragraph{HD-G00 (the role of Corollary O$''$).}
For $n=4$, the supplied Corollary O$''$ of
the supplied ATTACK gives $m\ge(16-8+3)/2$, hence $m\ge6$.
At $m=6$ it permits at most one nonreflection. The parabolic
dimension formula for an irreducible genus-zero local system is
$\dim V=\sum_i\operatorname{codim}\operatorname{Fix}(g_i)-2m$;
twenty reflections therefore give $\dim_KV=8$ and a Hurwitz source
of dimension $20-3=17$. These are the numerical conditions used to
select the datum below. The exact differential rank is established
separately in HD-G03--HD-G04 by that datum.

\subsection{The finite datum and the Hodge numbers}

Put $K=\mathbb Q(\omega)$, where $\omega^2+\omega+1=0$, and equip
$K^6$ with $\langle x,y\rangle=\sum_i\bar x_i y_i$. Define
\[
\begin{array}{ll}
v_1=(1,-1-\omega,1,\omega,-1-\omega,\omega),&
v_2=(1,0,0,1+\omega,0,0),\\
v_3=(0,1,1+\omega,0,0,0),&
v_4=(1,\omega,-1-\omega,\omega,-1-\omega,1),\\
v_5=(1,-1-\omega,-1-\omega,-1-\omega,-1-\omega,\omega),&
v_6=(1,-1-\omega,\omega,\omega,\omega,\omega),\\
v_7=(0,0,1,0,0,-1),&
v_8=(1,\omega,\omega,-1-\omega,1,-1-\omega),\\
v_9=(1,-\omega,0,0,0,0),&v_{10}=(0,1,0,0,0,-1).
\end{array}
\]
Let $s_i=1-2v_i\bar v_i^{,t}/\langle v_i,v_i\rangle$ and
$G=\langle s_1,\ldots,s_{10}\rangle$.

\paragraph{HD-G01 (group identification; exact computation and proof).}
The group $G$ is the Shephard--Todd group $G_{34}$ in its irreducible
six-dimensional reflection representation. Its order is $39\,191\,040$,
its centre is $\mu_6$, and its character field is $K$.

\paragraph{Proof.}
Each $s_i$ is a unitary involution with negative eigenspace $Kv_i$.
The exact calculation finds that the orbit of $v_2$ has 756 vectors,
and that the faithful permutation group on this orbit has order
$39\,191\,040$. Faithfulness follows because the orbit spans $K^6$:
the matrix calculation proves rank six, and any linear map fixing this
orbit fixes its linear span. The order is obtained by the ordinary
Schreier--Sims algorithm in SymPy; all arithmetic preceding this finite
permutation calculation uses fractions in $K$. It is an exact computer
certificate using rational arithmetic and the finite permutation algorithm.

For completeness, irreducibility follows from the checked rank-six span
of the ten roots and the connectedness of their nonorthogonality graph.
An invariant complex subspace $W$ has an invariant orthogonal complement.
The projector $(1-s_i)/2$ preserves both spaces and has one-dimensional
image, so each root line belongs to $W$ or to $W^\perp$. Nonorthogonal
root lines must belong to the same member of this decomposition. The
connected graph and spanning condition force $W=0$ or $W=\mathbb C^6$.

The published classification is Shephard--Todd, \emph{Finite unitary
reflection groups}, Canad. J. Math. 6 (1954), 274--304, \S\S2--4 and
Table VII, printed p.~301, row 34. In dimension six the other exceptional
group is $G_{35}$, of order $51\,840$, and the symmetric reflection group
has order $7!=5040$. An imprimitive group has order $720d^6/e$, with
$e\mid d$. Here $39\,191\,040/720=2^5 3^5 7$. Its 7-adic valuation is
one; the valuation of $d^6/e$ is either zero or at least five, because
$0\le v_7(e)\le v_7(d)$. This excludes every imprimitive candidate.
Thus the classified group is $G_{34}$.

The certificate tests membership of the six scalar units
$1,-1,\omega,-\omega,\omega^2,-\omega^2$. They all belong to $G$.
Schur's lemma makes every central element scalar, and a scalar finite-order
matrix over $K$ belongs to $\mu_6$. This proves the centre assertion.
All character values lie in $K$, and the trace of $\omega I_6$ is
$6\omega$, proving equality of the character field with $K$.

\paragraph{HD-G02 (the twenty-branch source).}
The tuple
\[
\boldsymbol g=(s_1,s_1,s_2,s_2,\ldots,s_{10},s_{10})
\]
defines connected $G_{34}$-covers of $\mathbb P^1$ branched at twenty
distinct points. Their Hurwitz component has dimension 17. The associated
primitive matrix-idempotent factor of the Jacobian has dimension eight
and $K$-signature $(4,4)$.

\paragraph{Proof.}
The tuple generates $G$ and has product one. Riemann existence therefore
constructs the covers (SGA 1, Expose XII, Theorem 5.1, applied off the
branch points, followed by normalization). The choices over an ordered configuration of twenty
points are finite; after marking a point above a fixed base point they
form a finite covering locally. Modulo $\operatorname{PGL}_2$, the
configuration dimension is $20-3=17$.
The reflection representation $\rho$ is nontrivial and irreducible.
Every local monodromy has eigenvalues $(-1,1,1,1,1,1)$, so its contribution
to the Chevalley--Weil holomorphic multiplicity is $1/2$.
The multiplicity is consequently $-6+20/2=4$ for both $\rho$ and
$\bar\rho$. This is the published Chevalley--Weil formula in
Frediani--Ghigi--Penegini, \emph{Shimura varieties in the Torelli locus via
Galois coverings}, IMRN 2015, Theorem 2.10; the preprint's printed p.~8
contains the same numbered formula.

The rational group-algebra factor of the pair $\rho,\bar\rho$ is
$M_6(K)$: the displayed representation is defined over its character
field and is absolutely irreducible. Let $\varepsilon$ be its rational
central idempotent and set $e=\varepsilon(1-s_1)/2$. Its action on
$K^6$ is the orthogonal projection onto $Kv_1$. The abelian subvariety
$A=\operatorname{im}(|G|e:J(C)\to J(C))$ therefore has first homology
equal to the multiplicity space, of $K$-dimension eight. The two
holomorphic multiplicities just computed give signature $(4,4)$.
The restricted theta polarization has Rosati involution equal to complex
conjugation on $K$, since $e$ is self-adjoint and the group acts by
polarization-preserving automorphisms. The central scalar $\omega$ gives
the integral action of $\mathbb Z[\omega]$ on $A$.

\subsection{Exact differential rank}

\paragraph{HD-G03 (parabolic monodromy certificate).}
In the order shown, take the following intervals of the twenty punctures:
\[
\begin{split}
&(1,8),(1,18),(2,8),(3,8),(4,8),(5,6),(5,8),(7,8),\\
&(9,10),(9,12),(9,16),(9,17),(9,18),(11,12),(13,14),(13,16),(15,16).
\end{split}
\]
Their boundary curves are disjoint or nested and form a maximal system
of seventeen essential curves. Let $N_j$ be the logarithm of the first
return of the corresponding full twist on the parabolic multiplicity
space $V$. Then
\[
\dim_KV=8,\quad N_jN_k=0,\quad
\dim_K\sum_j\operatorname{im}N_j=4,\quad
\operatorname{rank}_K\sum_jN_j=4,
\]
and
\[
\dim_K\langle N_1,\ldots,N_{17}\rangle_K
=\dim_{\mathbb Q}\langle N_1,\ldots,N_{17}\rangle_{\mathbb Q}=16.
\]
The individual ranks are
\[
(4,1,3,3,2,1,2,1,1,2,4,4,4,1,1,2,1).
\]

\paragraph{Proof and fully specified algorithm.}
Write $u_i=x_iv_{\lceil i/2\rceil}$ for the value on the $i$th standard
loop of a parabolic cocycle for the rank-six local system with monodromy
$\boldsymbol g$. Put $P_i=g_1\cdots g_{i-1}$. The relation matrix is
\[
C=(P_1v_1,P_2v_1,P_3v_2,P_4v_2,\ldots,P_{20}v_{10}),
\]
and the scalar coboundary matrix has row
$-2\bar v_{\lceil i/2\rceil}^{,t}/\langle v_{\lceil i/2\rceil},
v_{\lceil i/2\rceil}\rangle$.
Thus $V=\ker C/\operatorname{im} B$.
The calculation proves $CB=0$, $\operatorname{rank}C=6$,
$\operatorname{rank}B=6$; its quotient has dimension $20-6-6=8$.
This cocycle realization is also Dettweiler--Wewers,
\emph{Variation of local systems and parabolic cohomology}, Israel J.
Math. 156 (2006), Lemma 1.2 and the quotient immediately following it,
and Lemma 2.2/Theorem 2.5 for monodromy.
The comparison with the Jacobian factor is obtained by applying the
exact functor $\operatorname{Hom}_G(\rho^\vee,-)$ to $\pi_*K$ for
$\pi:C\to\mathbb P^1$. Finiteness gives $R\pi_*K=\pi_*K$, so the
cohomology of the resulting direct-summand sheaf is the corresponding
multiplicity space in $H^1(C,K)$. On the complement of the branch
points, $\pi_*K$ is the regular representation sheaf. Its matrix
coefficient decomposition identifies this summand with the rank-six
local system used above, up to simultaneous contragredience according
to the chosen regular-action convention. For these unitary matrices,
$\rho^\vee(g)=\overline{\rho(g)}$; the explicit field automorphism
$\omega\mapsto\omega^2$ identifies the conjugate calculation and
preserves every rank and product asserted here. Both conjugate
representations belong to the same rational idempotent $\varepsilon$.
This gives the rational multiplicity-space comparison used below.

Here is the precise operator used in the calculation. For an interval
$[a,b]$, put
\[
c=g_a\cdots g_b,\qquad
H(x)=\sum_{i=a}^b g_a\cdots g_{i-1}u_i.
\]
The first return exponent is $k=1$ for fourteen intervals and $k=2$ for
$[2,8],[4,8],[9,17]$. In the latter cases $c^2=1$ and at least one
$g_i$ in the interval does not commute with $c$, as checked exactly.
The $k$th twist pulls a cocycle back by
\[
u_i\longmapsto u_i+(1-g_i)(1+c+\cdots+c^{k-1})H(x)
\quad(a\le i\le b),
\]
and fixes the other coordinates. This follows directly from
$u(xy)=u(x)+\rho(x)u(y)$ applied to
$\gamma^k\ell_i\gamma^{-k}$, with
$\gamma=\ell_a\cdots\ell_b$ and $c^k=1$.

The script verifies preservation of the relation and pointwise fixation
of the coboundaries, then induces this operator on an elimination-chosen
quotient basis. Its difference from the identity squares to zero,
so that difference is the logarithm. All pairwise products, all ranks,
and both field-span ranks are then computed by exact elimination.
The independent logarithms have indices
\[
1,2,3,4,5,6,7,9,10,11,12,13,14,15,16,17.
\]
The basis and all seventeen matrices are included in the JSON certificate
as rational coefficient pairs in $1,\omega$. This establishes every
displayed finite assertion without a numerical tolerance.
There is a small exact rank witness: flatten each of the sixteen listed
logarithms, keep the entries in rows $1,2,3,4$ and columns $5,6,7,8$
in row-major order, and form the resulting $16\times16$ matrix.
Its determinant in this independently chosen quotient basis is $1728$.

\paragraph{HD-G04 (period differential has rank 16).}
The eightfold period map of HD-G02 has differential rank exactly 16 on
a nonempty open subset of its Hurwitz component. Its rational Hermitian
space is split.

\paragraph{Proof.}
Pinch the seventeen curves of HD-G03. Because there are $20-3$ disjoint
essential curves, this gives a stable rational curve with trivalent
components and seventeen independent node-smoothing parameters $q_j$.
Their existence follows from the deformation of a node $xy=q_j$;
the standard moduli statement is Knudsen, \emph{The projectivity of the
moduli space of stable curves II}, Math. Scand. 52 (1983), Theorem 2.7,
printed p.~179.
The Hurwitz family over the punctured polydisc becomes a family after
a finite Kummer cover. The local monodromy acts on parabolic cocycles
by the twist just computed, up to a finite scalar coming from the
central deck group. Marking a point above a base point removes that
scalar, or a further finite power kills it. Consequently the geometric
logarithms on this cover are $M_j=\epsilon c_jN_j$, with $c_j$ positive
integers and one common sign $\epsilon$. To see the common sign,
use the local cluster coordinate $z=q_jw$: a counterclockwise turn
of each $q_j$ rotates its entire cluster counterclockwise, producing
the same orientation of full twist for all nodes. Nested cluster
coordinates are products of these same $q_j$ and have the same property.

Let $I=\sum_j\operatorname{im}N_j$. It is four-dimensional, every
$N_j$ kills it, and $\sum N_j$ has kernel and image equal to $I$.
Choose positive integers $a_j$ with $a_jc_j$ independent of $j$ and
restrict the family to $u_j=t^{a_j}$ in the Kummer coordinates.
The resulting one-variable logarithm is a nonzero common scalar
multiple of $\sum N_j$, so its weight filtration, centred at one, is
\[
W_0=W_1=I,\qquad W_2=V.
\]
The limiting Hodge filtration $F_0$ of the several-variable map
is also the limiting filtration on this diagonal curve. By Schmid,
\emph{Variation of Hodge structure: the singularities of the period
mapping}, Invent. Math. 22 (1973), Theorem (4.12), printed p.~233,
the untwisted map extends holomorphically; by Theorem (6.16),
printed p.~255, $(W,F_0)$ is a mixed Hodge structure.
These statements are applied to the underlying rational first cohomology;
we then take its chosen $K$-eigenspace.
Since $F_0^0=V_{\mathbb C}$ and $F_0^2=0$, its weight-zero graded piece
has only type $(0,0)$ and its weight-two graded piece only type $(1,1)$.
On the chosen $K$-eigenspace this gives
\[
F_0^1\cap I_{\mathbb C}=0,\qquad
\dim F_0^1=4,\qquad V_{\mathbb C}=I_{\mathbb C}\oplus F_0^1.
\]

The matrices $N_j$ map $V/I$ to $I$ and span a 16-dimensional
$K$-space. Since $\operatorname{Hom}_K(V/I,I)$ has dimension $4\cdot4=16$,
they span that entire space. The explicit isomorphism to the tangent
space at $F_0^1$ is
\[
\operatorname{Hom}_{\mathbb C}(V_{\mathbb C}/I_{\mathbb C},I_{\mathbb C})
\longrightarrow
\operatorname{Hom}_{\mathbb C}(F_0^1,V_{\mathbb C}/F_0^1),
\quad X\longmapsto \bigl(F_0^1\hookrightarrow V/I
\xrightarrow{X}I\hookrightarrow V/F_0^1\bigr).
\]
Both outside arrows are isomorphisms by the displayed direct sum.
Schmid's theorem writes the lifted map as
\[
\Phi(z)=\exp\!\left(\sum_j z_jM_j\right)\psi(u),
\qquad u_j=e^{2\pi iz_j},\quad\psi(0)=F_0.
\]
After translating tangent vectors by the inverse exponential, the
$j$th derivative is
$M_j\psi(u)+2\pi i u_j\partial\psi/\partial u_j$.
It tends to $M_jF_0$ as $u\to0$. Sixteen columns have a nonzero
limiting determinant; their determinant remains nonzero at nearby
interior points. The target Weil domain has dimension 16, proving
the exact rank.

Each geometric monodromy preserves the polarization; its logarithm is
skew-adjoint for the $K$-Hermitian form $h$ on $V$. Consequently
\[
h(N_jx,N_ky)=-h(x,N_jN_ky)=0.
\]
The explicit subspace $I$ is therefore a rational totally isotropic
$K$-subspace of half the dimension. This proves splitness. The
identification of $V$ with the Jacobian multiplicity space respects
the polarization up to a nonzero rational scalar and, according to
the dual convention, conjugation of $K$; both operations preserve the
existence of this rational Lagrangian. Thus the eightfold Hermitian
space is split.

\subsection{Cycles without an unproved supplier premise}

\paragraph{HD-G05 (exceptional-divisor cohomology in the supplier construction).}
For every cyclic intermediate quotient of the $G$-covers in HD-G02,
the primitive-character vanishing required in Schoen's Lemma 2.8 holds,
including ramified quotients of cyclic order divisible by four.

\paragraph{Proof.}
Let $H\le G$ and let $\chi:H\to\mu_m$ have exact order $m\ge3$.
Put $D=C/\ker\chi$ and $X=C/H$. Every inertia group of $C\to\mathbb P^1$
has order two. Therefore every nontrivial inertia element of $D\to X$
has order two and exponent $m/2$; when $m$ is odd the cover is unramified.
The product relation in the abelian group $\mu_m$ makes the number
of nontrivial exponents even. Thus the exponents pair with their negatives.
The primitive multiplicity is $h=2g(X)-2+s$, an even integer, where $s$
is the number of branch points. Terms with $h=0$ will be discarded below.

Use Schoen, Compositio Math. 65 (1988), Theorem 2.0, Lemmas 2.3--2.6,
printed pp.~11--17, to construct the linear space and its split lifts.
At a blow-up of the intersection of the two branch divisors in one
opposite-exponent pair, let $p:E\to Z$ be the exceptional $\mathbb P^1$
bundle and $F$ its inverse image in the normalized cyclic cover.
The two exponents sum to zero, so $E$ is generically unramified
(Schoen, Lemma 2.5). The branch divisors are a normal-crossing arrangement:
the paired divisors cut out two disjoint sections in $E$, and every
other branch divisor meets $E$ in the inverse image of a closed subset
of $Z$. This follows in the normal-crossing coordinates by blowing up
the ideal $(x_1,x_2)$; a remaining divisor has local equation $x_j=0$
with $j>2$, independent of the exceptional direction.

For a primitive character $\eta$ of $\mu_m$, let $\mathcal F$ denote
the $\eta$-eigensheaf in the finite direct image of the constant sheaf
from $F$ to $E$. At a branched point its stalk is zero: the fibre is a
permutation representation induced from its inertia subgroup, and a
faithful character has no invariants on a nontrivial inertia subgroup.
On the remaining open subset it is a rank-one local system, extended
by zero across the branch locus.
On a fibre of $p$ contained in an additional branch divisor its
restriction is zero. On any other fibre it is the extension by zero
of a rank-one local system on $\mathbb P^1\setminus\{0,\infty\}$ with
monodromy $\eta(m/2)=-1$. Its ordinary cohomology is computed by
the two-term cellular complex $\mathbb C\xrightarrow{1-(-1)}\mathbb C$,
which is acyclic. Poincare duality gives zero compactly supported
cohomology as well.
Proper base change therefore gives $Rp_*\mathcal F=0$, and hence
$H^*(F,\mathbb C)_\eta=H^*(E,\mathcal F)=0$.
The precise general topological input is SGA 4, Expose V bis,
4.1.1 (also Stacks Project, Theorem 20.18.2, tag 09V6);
the etale counterpart is SGA 4, Expose XII, Theorem 5.1.
This proves the vanishing that Schoen uses to descend the nonzero
primitive cycle classes through his blow-ups. Together with his
published normal-bundle computation it supplies those classes for
every cyclic quotient used here. In particular no selection of a
prime quotient ramified at the branch point is required.

\paragraph{HD-G06 (Brauer transfer for this source).}
For every cover of HD-G02, the Weil plane of its eightfold factor is
spanned by algebraic cycle classes.

\paragraph{Proof.}
Let $V_C=H^1(C,\mathbb C)$, with its cup-product alternating form,
and let $U$ be the centralizer of $G$ in its symplectic group.
For a finite-dimensional complex $G$-module $R$, put
$M_R=\operatorname{Hom}_G(R,V_C)$. The assignment
$R\mapsto\det M_R$ is an additive homomorphism from the representation
ring to the character group of $U$: direct sums in $R$ become direct
sums in $M_R$, and determinants become products. Brauer's induction
theorem (Serre, \emph{Linear representations of finite groups},
\S10.5, Theorem 20, printed pp.~78--79) expresses the reflection representation as
\[
[\rho]=\sum_i n_i[\operatorname{Ind}_{H_i}^G\chi_i],\qquad n_i\in\mathbb Z,
\]
with linear characters $\chi_i$. Frobenius reciprocity then identifies
the determinant character of $M_\rho$ with the product of the
characters $\det\operatorname{Hom}_{H_i}(\chi_i,V_C)$ to the powers $n_i$.
For $\chi_i$ of order one or two, that multiplicity space has a
nondegenerate alternating form obtained by restriction of the cup
product, and $U$ acts symplectically; its determinant character is
one. A zero multiplicity space has determinant character one as well.
Discard these terms. Each remaining term is a positive-dimensional
cyclic supplier covered by HD-G05.

Write $P_i\subset J(C)$ for its primitive abelian factor. Its generalized
Weil space is a one-dimensional vector space over $\mathbb Q(\mu_{m_i})$.
Schoen's Theorem 2.0 and Corollary 3.1, with HD-G05 at the exceptional
step, give a nonzero rational algebraic class $y_i$ in that space.
Its complexification has nonzero component in every primitive eigenline:
a nonzero vector in a line over a number field stays nonzero under each
embedding. Choose the eigenline $\chi_i$ for $n_i>0$ and its conjugate
for $n_i<0$, and let $Z_i$ be the wedge product of all the other
primitive determinant lines in $H^*(P_i,\mathbb C)$. The full product
is top cohomology, whose $U$-character is trivial. Thus $Z_i$ has the
inverse character to the chosen line and pairs nontrivially with the
chosen component of $y_i$; it annihilates the other components by
exterior-degree repetition.

Let $B=\prod_iP_i^{|n_i|}$, let $Y=\boxtimes_i y_i^{\boxtimes|n_i|}$,
and let $x$ span $\det M_\rho$, one complex Weil line of $A$. Set
\[
\Gamma=\left(\boxtimes_iZ_i^{\boxtimes|n_i|}\right)\boxtimes x.
\]
The determinant-character identity makes $\Gamma$ invariant under $U$.
Every factor of $B\times A$ is the image of a rational idempotent
in $\mathbb Q[G]$. Choose a common integer clearing these idempotents;
the inclusions into and quotient maps from a power $J(C)^r$ are
$U$-equivariant and their composite is multiplication by that integer.
Pulling $\Gamma$ back to $J(C)^r$ gives a $U$-invariant class.
Milne's group $S(J(C))$, the subgroup of the symplectic group commuting
with every endomorphism of $J(C)$, is contained in $U$. Milne,
\emph{Lefschetz classes on abelian varieties}, Duke Math. J. 96 (1999),
Lemma 3.1 and Theorem 3.2, therefore express the pullback as a complex
linear combination of products of divisor classes. Pullback along the
inclusions recovers a nonzero integer scalar multiple of $\Gamma$,
so $\Gamma$ is algebraic with complex coefficients.

The correspondence action of $\Gamma$ on $Y$ is a nonzero scalar multiple
of $x$, by the stated eigenline pairings and the projection formula.
Thus $x$ is algebraic with complex coefficients. Its complex conjugate
is algebraic as well. If $\mathcal Z_{\mathbb Q}$ is the rational vector
space of algebraic classes, this proves
$W_K(A)\otimes\mathbb C\subseteq\mathcal Z_{\mathbb Q}\otimes\mathbb C$;
faithful scalar extension yields $W_K(A)\subseteq\mathcal Z_{\mathbb Q}$.

\subsection{An integral uniform detector bound}

\paragraph{HD-G07 (checked quantitative extraction).}
Let $G=39\,191\,040$ and $g=4G+1$. Define integers
\[
R=(2G+1)G!G^G,\qquad L=2gR,\qquad
S_0=G^{2g}(g!)^2,\qquad S_B=(L/2)!S_0^R,
\]
\[
k=L/2+4,\qquad
B_0=(40G)^k\binom{k}{4}S_B\,8!G^8.
\]
Every eightfold $A$ of HD-G02 has an effective integral four-cycle
with nonzero Weil projection and degree at most $B_0$ for the restricted
Jacobian theta polarization.

\paragraph{Proof and provenance.}
The quantitative argument is HECKE\_DIV SQ2--SQ6, whose proofs were
read in full and rechecked in this run; their two previously unrefereed
dependencies are now HD-G05 and HD-G06. We record the exact estimates
to specify their use. The self-adjoint idempotent satisfies $Ge\in\mathbb Z[G]$,
and the polarization elementary divisors of $A$ divide $G$; hence
$\int_A\theta_A^8\le8!G^8$ (SQ2). Integer linear algebra applied to
the Brauer character lattice gives a supplier identity with total
absolute coefficient at most $R$ (SQ3). For a supplier of order $m$
and multiplicity $h$, SQ4 constructs an integral signed cycle on $C'^h$
of degree mass at most $m^{3h/2}(h!)^2\le S_0$ for the sum of point classes.
Its exterior products use at most $L$ curve factors and have mass
at most $S_B$ (SQ5).

The invariant transfer tensor is spanned by products of globally
generated divisor classes pulled back from $2\Theta$ on $J(C)$ along
the maps $x_i$ and $x_i+ax_j$ for $a\in G$, including repeated indices.
Milne's Lemma 3.5, Proposition 3.6 and Remark 3.7 give this invariant
theory statement. On the product of the curve supplier and $A$, each
such divisor is at most $40G(H_B+\theta_A)$ in the nef order: a quotient
curve has theta degree at most $5G$, and the theorem of the square
gives the factor eight. Proper intersections with general members
of these globally generated systems and pushforward to $A$ are
effective integral cycles. For $c=\ell/2+4\le k$, the mass is at most
\[
(40G)^c\binom c4 S_B\int_A\theta_A^8\le B_0.
\]
The transfer sends the supplier class to a nonzero Weil class, so some
effective component in this finite spanning construction has nonzero
Weil projection. This is precisely the SQ6 nonvanishing argument;
the published suppliers and the proved HD-G06 establish its hypotheses.

\paragraph{HD-G08 (uniform bounded division for split Eisenstein eightfolds).}
Fix a polarized split Weil eightfold family with integral
$\mathbb Z[\omega]$-action and polarization elementary-divisor product
$\Delta$. Put
\[
D_8(\Delta)=(4G^8\Delta)^8 B_0.
\]
Then $\delta_W\le D_8(\Delta)$ on every fibre of the family. For a fixed
nonzero rational Weil class $t$, with $\delta t$ integral and
$T=\int t^2>0$, Theorem W of the supplied ATTACK R4 gives the
integral effective bound
\[
c_t\le
\left\lceil C_{4,3+\omega}(\delta,T)D_8(\Delta)^4\right\rceil.
\]

\paragraph{Proof.}
HECKE\_DIV SQ7 gives an explicit rational Hermitian isometry from
the source lattice to a fixed target lattice, cleared in both directions
by $N=4G^8\Delta$. Its proof uses the matrix
$\begin{pmatrix}0&B\\B^*&D\end{pmatrix}$ in a primitive isotropic
$\mathbb Z[\omega]$-basis and the substitutions
$w'=w-v(B^*)^{-1}D/2$, $u=w'B^{-1}$; both coordinate maps are
cleared by $2\Delta$. The source elementary-divisor product divides
$G^8$, so $N$ clears their comparison. This entire lattice calculation
was rechecked, including the normalization
$h(x,y)=E(x,\omega y)-\bar\omega E(x,y)$ and
$|\det h|=\Delta$.
On a simply connected local source, the resulting isogeny satisfies
$\phi^*\theta_{\rm target}=N^2\theta_A$.
Consequently an effective four-cycle pushes forward with degree
multiplied by $N^8$, and its Weil projection remains nonzero.
HD-G04 gives a nonempty open image in the target Weil domain.
Apply the next step on its connected algebraic moduli component, after
finite level, and then restrict to the stated family. The fixed rational
comparison is realized by a bounded-degree isogeny correspondence; hence
the analytic open image lies in that moduli component and is Zariski dense.
Thus HD-G07 yields the displayed detector bound on a dense subset of
that component.
The detector closedness and Theorem W are the proved inputs
the supplied ATTACK Theorem D and R4, so they extend this bound
to the whole family and produce the stated division bound.

\paragraph{HD-G09 (a fixed nonmaximal order).}
Let the split Eisenstein eightfold family carry the integral action of
$\mathcal O_f=\mathbb Z+f\mathbb Z[\omega]$, with conductor $f\ge1$,
and let its polarization elementary-divisor product be $\Delta$.
Then every fibre has
\[
\delta_W\le f^{60}D_8(\Delta),\qquad
c_t\le\left\lceil
C_{4,f(3+\omega)}(\delta,T)
\bigl(f^{60}D_8(\Delta)\bigr)^4\right\rceil.
\]

\paragraph{Proof.}
Write its lattice as $\Lambda$, with integral alternating form $E$, and put
$\Lambda'=\mathbb Z[\omega]\Lambda=\Lambda+\omega\Lambda$.
Since $f\omega$ acts integrally, $f\Lambda'\subseteq\Lambda$.
Thus $\Lambda'$ defines an algebraic quotient isogeny
$q:A\to A'$ with finite degree $d=[\Lambda':\Lambda]$;
its kernel is the displayed subgroup of $A[f]$.
On $\Lambda'$ use $E'=fE$. This is integral: the mixed terms
$fE(x,\omega y)$ and $fE(\omega x,y)$ are integral by the
$\mathcal O_f$-action, and
$E(\omega x,\omega y)=E(x,y)$ by Rosati compatibility.
It is a positive compatible polarization, denoted $\theta'$, and
$q^*\theta'=f\theta$. The target has integral maximal-order action
and the same split rational Hermitian space. Its polarization product is
\[
\Delta'=\frac{f^8\Delta}{d},
\]
as follows by integrating $(q^*\theta')^8$ and dividing by $8!$.
Take the effective integral detector $Z'$ furnished by HD-G08.
Finite flat pullback is effective and integral and preserves its nonzero
Weil projection. Its degree satisfies
\[
\deg_\theta q^*Z'
=\frac d{f^4}\deg_{\theta'}Z'
\le\frac d{f^4}D_8(\Delta')
=\frac{f^{60}}{d^7}D_8(\Delta)
\le f^{60}D_8(\Delta).
\]
The equality uses the explicit homogeneity $D_8(x)=\text{constant}\cdot x^8$.
An integral component of this pullback is a detector of no larger degree.
Finally $f(3+\omega)\in\mathcal O_f$ has quotient by its conjugate
$(8+5\omega)/7$. This is outside $\mathbb Z[\omega]$, whereas every
root of unity in $K$ lies in that ring. Thus the quotient is not a
root of unity, and the stated form of Theorem W applies.

\paragraph{Framework location.}
The Hurwitz construction consists of finite covers and idempotent
summands of a relative Jacobian. Its homological motive is cut out by
$e$. On the rank-eight $K$-local system the determinant is
$\lambda_K^8(V)=\bigwedge_K^8V$; the Brauer identity of HD-G06 is an
identity in $K_0(\operatorname{Rep}_{\mathbb C}G)$ followed by its
explicit determinant-character homomorphism.
The realization in curve powers has a precise sign convention:
on the $H^1$ tensor piece of $C^r$, a geometric permutation $\pi$
acts as $\operatorname{sgn}(\pi)$ times the ordinary permutation of
tensor factors. Hence the geometric symmetrizer
$r!^{-1}\sum_{\pi\in S_r}\pi^*$ realizes the ordinary exterior-power
projector there. This specifies the motivic projector producing the
determinant summands.
The suppliers and correspondence $\Gamma$ are
algebraic cycle morphisms of these motives. Finite-type descent spreads
each finite construction to a localization of a finitely generated
$\mathbb Z$-algebra (EGA IV$_3$, Theorem 8.8.2(i)--(ii), descent of
finitely presented schemes and their morphisms); the exceptional-divisor
vanishing uses proper base change (SGA 4 XII 5.1, with the Betti form
above). The parabolic local system is the corresponding summand of
$R^1$ of the cover family; its geometric fibres have weight one, and
its determinant realization has weight eight, or weight zero after
the twist $(4)$. At a good finite-field fibre its determinant Euler
factor, on a chosen $K$-embedding, is
$1-\det(\operatorname{Frob}\mid V_\ell)T$.
The argument constructs
the cycles used for Grothendieck's condition A; no estimate of an
$L$-value supplies a cycle here. In the project's archive frame, this
is cote 16, scan 050, condition A as corrected in SIGN\_NEXT SN-C03/C04.
The subsequent positive Gram-matrix division is condition I, cote 16,
scan 049, and cote 44, Theorem 4. Theorem W's resultant projector is
the explicit correspondence-ring inverse associated with the
Cayley--Hamilton construction in cote 16, scans 048--049.
These archive identifications are carried from the supplied corrected
intakes; no original archive scan was supplied or independently read
in this subtask.
